\section{Learning Local Policies in Markov Bidding Games.}
Our goal is to learn the local policy of each player in the constructed Markov game using RL.
\citet{yu2022mappo} demonstrate that proximal policy optimization (PPO; \cite{schulman2017ppo}) can be effective to train multiple policies concurrently for multi-agent settings. We follow their recipe and train all the local policies with PPO. 
Similar to \citet{yu2022mappo}, since the objectives in our environments share the same characteristics, we use the same actor and critic network for all the policies. 
At each iteration of PPO, we collect all the rollouts from each policy's perspective and sample minibatches from this pool to perform the PPO updates.

Additionally, to cope with varying number of objectives and varying number of players in the resulting Markov game, we employ the \textit{attention pooling} architecture to compress any number of objectives into a fixed-sized embedding. 
This design enables the agent to operate with an arbitrary number of policies (i.e., objectives) at deployment time. 
Each objective information vector is first processed by a fully connected neural network to produce a target embedding. 
These embeddings are then aggregated through a weighted sum, where the weights are computed by taking the dot product between each embedding and a learned query vector. 
This attention mechanism produces a fixed-dimensional representation that summarizes the set of objectives while allowing the model to adapt to varying objective counts.
\GS{Talk about discrete bids.}
